\documentclass[11pt,letterpaper]{article}
\usepackage[margin=0.85in,headheight=16pt]{geometry}
\usepackage[T1]{fontenc}
\usepackage[utf8]{inputenc}
\usepackage{lmodern,microtype}
\usepackage{amsmath,amssymb,mathtools,booktabs,longtable,tabularx,array}
\usepackage[dvipsnames]{xcolor}
\usepackage{hyperref,fancyhdr,titlesec,enumitem,graphicx,tcolorbox}
\usepackage{lastpage}
\definecolor{ink}{HTML}{14263B}
\definecolor{sea}{HTML}{167C8A}
\definecolor{gold}{HTML}{C49747}
\definecolor{pale}{HTML}{EEF5F6}
\definecolor{graytext}{HTML}{516275}
\hypersetup{colorlinks=true,linkcolor=ink,urlcolor=sea,bookmarksopen=true,pdfauthor={Independent Educational Engineering Suite}}
\titleformat{\section}{\Large\bfseries\color{ink}}{\thesection}{0.7em}{}[\vspace{-0.3em}\color{gold}\titlerule]
\titleformat{\subsection}{\large\bfseries\color{sea}}{\thesubsection}{0.6em}{}
\setlist[itemize]{leftmargin=1.3em,itemsep=2pt,topsep=3pt}
\setlist[enumerate]{leftmargin=1.5em,itemsep=2pt,topsep=3pt}
\pagestyle{fancy}
\fancyhf{}
\fancyhead[L]{\small\sffamily\color{graytext}SECURITY ANALYSIS / ENGINEERING SUITE}
\fancyhead[R]{\small\sffamily\color{graytext}\textbf{\thepage}/\pageref{LastPage}}
\renewcommand{\headrulewidth}{0.3pt}
\renewcommand{\footrulewidth}{0pt}
\newtcolorbox{methodbox}[1][]{colback=pale,colframe=sea,boxrule=0.55pt,arc=1mm,title={\bfseries Method note},fonttitle=\sffamily,#1}
\newtcolorbox{sourcebox}[1][]{colback=white,colframe=gold,boxrule=0.55pt,arc=1mm,title={\bfseries Source boundary},fonttitle=\sffamily,#1}
\newcommand{\term}[1]{\textbf{\color{ink}#1}}
\newcommand{\src}{\textsf{\color{sea}[SOURCE]}}
\newcommand{\ext}{\textsf{\color{gold}[EXTENSION]}}
\newcommand{\gap}{\textsf{\color{red}[SOURCE GAP]}}
\newcommand{\answer}[1]{\par\smallskip\noindent\textbf{Check:} #1\par}
\newcommand{\cover}[3]{%
\begin{titlepage}\pagecolor{ink}\color{white}\vspace*{1.15in}
{\sffamily\large\color{gold}INDEPENDENT TECHNICAL DOCUMENTATION SUITE\par}\vspace{0.4in}
{\Huge\sffamily\bfseries #1\par}\vspace{0.15in}
{\Large\sffamily #2\par}\vspace{0.65in}
\color{gold}\rule{\textwidth}{2pt}\par\vspace{0.4in}
{\large #3\par}\vfill
{\small\sffamily Based on a critical reading of Graham and Dodd's \emph{Security Analysis}, Seventh Edition (2023). Original educational analysis; not affiliated with the publisher.\par}
\end{titlepage}\nopagecolor\color{black}}
\newcommand{\front}[3]{\cover{#1}{#2}{#3}\tableofcontents\clearpage
\begin{sourcebox}\src\ means a concept is attributable to the relevant chapter or seventh-edition contributor, summarized independently rather than reproduced. \ext\ means a modern method or engineering recommendation prepared for this suite, not an original Graham--Dodd rule. Historical examples and standards should not be treated as current market requirements. All numerical examples are synthetic, educational, and not securities recommendations.\end{sourcebox}}
\setlength{\parskip}{5pt}\setlength{\parindent}{0pt}
\begin{document}
\front{Security Analysis}{Prerequisite-Based Technical Curriculum}{Sixteen-week progression, gates, laboratories, mastery rubric and capstone}

\section{Program overview and entrance requirements}
Audience: financial analysts, value-investing researchers, credit analysts, investment risk engineers, and advanced independent learners. Prerequisites: accounting statement literacy, algebra, time value of money, spreadsheet or scripting competency, and basic probability. Learners without prerequisites take Foundation F0 before Core Modules M1--M9.

\section{Prerequisite dependency architecture}
\begin{center}\fbox{\begin{minipage}{0.9\textwidth}\centering
\textbf{F0: Accounting, mathematics, spreadsheet controls}\\[4pt]
$\Downarrow$\\[-2pt]
\textbf{M1: Evidence + claims} $\longrightarrow$ \textbf{M2: Credit underwriting}\\[4pt]
$\Downarrow$ \hspace{2cm} $\Downarrow$\\[-2pt]
\textbf{M3: Earnings quality} $\longrightarrow$ \textbf{M4: Valuation models}\\[4pt]
$\Downarrow$ \hspace{2cm} $\Downarrow$\\[-2pt]
\textbf{M5: Asset realization} $\longrightarrow$ \textbf{M6: Distress + optionality}\\[4pt]
$\Downarrow$\\[-2pt]
\textbf{M7: Due diligence} $\longrightarrow$ \textbf{M8: Portfolio risk}\\[4pt]
$\Downarrow$\\[-2pt]
\textbf{M9: Integrated investment committee capstone}
\end{minipage}}\end{center}
The diagram shows primary dependencies. Credit, earnings and asset work can be studied concurrently after M1; M7 requires credible competency in M2--M6 before capstone underwriting.

\section{Suggested 16-week teaching schedule}
\begin{center}\small\begin{longtable}{p{0.11\textwidth}p{0.27\textwidth}p{0.23\textwidth}p{0.28\textwidth}}
\toprule Week & Module & Textbook map & Deliverable \\
\midrule\endhead
1 & F0: statements, discounting, model hygiene & Preface; Ch. 1--3 & Input validation workbook \\
2 & M1: claims and source reliability & Ch. 1--5 & Evidence and instrument map \\
3 & M2: bond selection and stress & Ch. 6--11 & Bond-underwriting screen \\
4 & M2: preferred and covenants & Ch. 12--21 & Covenant/capital stack file \\
5 & M6: convertible and distressed credit & Ch. 22--26 & Payoff and recovery case \\
6 & M3: equity logic and dividends & Ch. 27--30 & Earning-power thesis \\
7 & M3: income statement adjustments & Ch. 31--34 & Normalized earnings bridge \\
8 & M3: history, structure and dilution & Ch. 35--41 & Adjustments/dilution memo \\
9 & M5: balance sheet and NCAV & Ch. 42--45 & Asset realization sensitivity \\
10 & M4: DCF and peer selection & Ch. 39; 49 & Reproducible DCF/peer model \\
11 & M6: warrants, governance, mispricing & Ch. 46--52 & Three-scenario value-gap memo \\
12 & M7: business DD + source review & Ch. 3; 31--45 & Evidence registry and challenge log \\
13 & M8: portfolio and credit risk & Part II; Part VIII & Limit and stress dashboard \\
14 & M8: international/endowment extensions & Part VIII & Currency/liquidity stress \\
15 & M9: integrated capstone preparation & All assigned & Full research packet \\
16 & M9: committee defense and remediation & All assigned & Oral defense, revised thesis \\
\bottomrule\end{longtable}\end{center}
\emph{Instructor note:} Chapter 36 is listed in the supplied PDF contents but its text is absent. Teach only a clearly labeled independently developed amortization module unless the supplement is separately obtained and reviewed.

\section{Competencies and prerequisite checks}
\subsection{F0 --- foundational controls (6--10 hours)}
\textbf{Objectives:} reconstruct basic statements; discount dated cash flows; check financial units. \textbf{Gate:} 80\% on a 20-question accounting/calculation assessment plus error-free source-to-calculation tie-out.

\subsection{M1 --- instruments, evidence and investment logic (10 hours)}
\textbf{Objectives:} define intrinsic value and investor claim; classify investment/speculation; cite primary sources. \textbf{Gate:} independently identify three disclosure discrepancies and clearly explain the effect on a thesis.

\subsection{M2 --- fixed-income safety and covenants (16 hours)}
\textbf{Objectives:} evaluate trough coverage, collateral, parent/affiliate structure, preferred provisions and monitoring triggers. \textbf{Gate:} defend one credit under a 30\% earnings decline and an imminent maturity.

\subsection{M3 --- accounting and earnings quality (14 hours)}
\textbf{Objectives:} reconstruct operating earning power, separate exceptional items, explain working capital and dilution. \textbf{Gate:} produce a normalized three-year cash/earnings bridge reproducible to the dollar.

\subsection{M4 --- quantitative valuation (18 hours)}
\textbf{Objectives:} build FCFF DCF, interest-rate and terminal sensitivity, equity bridge, peer-normalized multiples and reverse DCF. \textbf{Gate:} explain how changing $r$, $g$, reinvestment and leverage changes the value range.

\subsection{M5 --- asset-based downside (10 hours)}
\textbf{Objectives:} compute NCAV, asset haircuts, disposal expenses and legal-entity claims. \textbf{Gate:} produce a pessimistic liquidation waterfall without omitting liabilities.

\subsection{M6 --- optionality, distress and price--value discrepancies (12 hours)}
\textbf{Objectives:} calculate parity, warrant dilution, simplified restructuring distribution, and justified versus unjustified discounts. \textbf{Gate:} distinguish creditor floor from contingent equity-option participation.

\subsection{M7 --- operational due diligence (10 hours)}
\textbf{Objectives:} maintain an evidence registry, adversarial review and investment memo. \textbf{Gate:} address at least five material counterarguments or record a justified rejection.

\subsection{M8 --- portfolio risk and supervision (12 hours)}
\textbf{Objectives:} calculate issuer exposure, drawdown, simple DV01, cash liquidity, and stress losses; evaluate model risk. \textbf{Gate:} perform a two-scenario stress and specify an implementable mitigation.

\subsection{M9 --- investment committee capstone (15--20 hours)}
\textbf{Objective:} compare one fictional equity issuer and one debt issuer in a constrained simulated portfolio. \textbf{Deliverables:} 12-page research memorandum, model with tests, claim waterfall, risk dashboard, 10-minute oral defense and revised thesis responding to dissent.

\section{Practical laboratories}
\begin{enumerate}
\item \textbf{L1 Evidence lab:} log four financial filings with consistent period keys and two independently confirmed adjustments.
\item \textbf{L2 Credit lab:} build three-year debt maturity ladder and calculate covenant headroom under two shock severities.
\item \textbf{L3 Earnings lab:} reconcile GAAP net earnings, normalized EBIT, recurring FCFF and fully diluted EPS.
\item \textbf{L4 Valuation lab:} compare discounted cash flow, multiples and liquidation support with uncertainty envelopes.
\item \textbf{L5 Optionality lab:} plot convertible parity against equity price and show downside credit sensitivity.
\item \textbf{L6 Portfolio lab:} combine correlated exposures; stress a redemption with liquidation haircuts.
\end{enumerate}

\section{Assessment and advancement policy}
\begin{center}\small\begin{tabularx}{\textwidth}{p{0.25\textwidth}p{0.18\textwidth}X}\toprule Criterion & Weight & Evidence \\
\midrule Accuracy and source integrity & 25\% & Tied-out primary input values \\
Valuation/recovery correctness & 25\% & Independently recomputable logic \\
Downside and uncertainty & 20\% & Credible loss, stress, sensitivity cases \\
Portfolio fit and governance & 15\% & Limits, liquidity and approvals \\
Written/oral argument & 15\% & Persuasive conclusion with dissent response \\
\bottomrule\end{tabularx}\end{center}
\textbf{Recommended pass standard:} total at least 80\%, no critical calculation or source-integrity failure, and all mandatory capstone documents submitted. Reassessment must show corrected model and cause of initial error, not merely a revised numerical answer.

\section{Instructor resources and progression}
Use this suite's \emph{Technical Reference} for source orientation, \emph{Chapter Study Guide} for reading checks, \emph{Quantitative Handbook} for formulas, \emph{Modeling Exercises} for graded labs, \emph{Due Diligence SOP} for capstone packet structure, and \emph{Portfolio Risk Framework} for committee challenge. After graduation, assign sector-specific credit, international governance, multi-asset liquidity or private investment electives.

\end{document}
